\chapter{Conclusions and Future Work}
\label{chap:conclusions}

This chapter closes the thesis: it checks the outcome against the
requirements of Chapter~\ref{chap:problem}, states the design lessons
the work argues for, and lays out the directions it opens.

\section{Conclusions}
\label{sec:conclusions}

This thesis set out from a one-sentence gap: the tools we surveyed answer
``what did my agent do?'' with a trace of one run on one machine,
and, to our knowledge, none answers ``what is my agent \emph{fleet}
doing across the cluster, and which node is dragging it down?'' The hypothesis defended across the preceding
chapters is that the gap closes naturally once the data is stored in its
native shape: a multi-agent execution is a causally-ordered event log, and
a distributed shared log with a capture process on every node makes the
coupling of agent semantics to physical topology a property
of the substrate rather than an integration problem.
Section~\ref{sec:problem-substrate} states the contract a substrate must
honour, which makes the hypothesis checkable against any candidate store, and
Chapter~\ref{chap:design} is written against it, with
\S\ref{sec:design-chronolog} isolating what implementing it on ChronoLog
in particular required.

All seven requirements of \S\ref{sec:problem-requirements} are met at the
scale evaluated, each
with a measured answer from Chapter~\ref{chap:evaluation}. \textbf{R1} is
met by the two capture paths and the \texttt{role@scenario} identity
contract, under which every event carries its host and the planes join by
construction. \textbf{R2} and \textbf{R3} are met jointly by the hot/cold
split: an event reaches the operator's screen 10.7\,ms after emission and
the durable archive 20--30\,s after, with at-least-once spooled capture
deduplicated by high-water marks rebuilt from the log itself. \textbf{R4}
required changing both computation and representation: per-(host, minute)
rollups hold the fleet read at 19--24\,ms across a $50\times$ growth in
event volume where the raw fold grows $457\times$ slower, and the heatmap
answers a 400-node fleet in 49\,ms where node-link rendering fails at a
dozen.
\textbf{R5} holds structurally and empirically: 3.6\,ms per
turn-plus-delegation, under 0.05\% of a mean turn, with an A/B delta
statistically indistinguishable from zero. \textbf{R6} rests on the
deterministic detect--cluster--diagnose pipeline, LLM assistance strictly
opt-in, scoring 580/580 planted faults with zero false positives and
compressing a 500-signal storm to one ranked card. \textbf{R7} is
demonstrated at that scale by the live harness: 28/28 automated checks on
4- and 8-node
clusters without modifying ChronoLog, the scheduler, or the agent SDK,
sustained at 32 concurrent real agents with flat per-node sidecar load.
Whether the same holds on an allocation an order of magnitude larger is
untested.

Beyond the artifact, the thesis argues two transferable design lessons.
The first is architectural. Observability systems are conventionally
assembled as instrumentation plus a bespoke backend; for workloads
whose events are few, expensive, and causally structured (as
LLM-agent fleets are), inverting the architecture around a general
distributed shared log yields the cross-layer joins, the durability,
and the replayability for free, at the price of engineering around the
log's read behavior. That price is real but modest, and the evaluation
quantified both of its parts: the planned part (the hot/cold split of
Chapter~\ref{chap:design}) and the unplanned part
(the replay behavior reported in \S\ref{sec:eval-limitations}, absorbed by
generalizing archive-first reads at a bounded cost of half a minute of
staleness). In return it provides the capability neither surveyed camp
offers: the pivot from a fleet-wide symptom to a single agent's
conversation on a named machine. The case study measured that pivot at
72 seconds, questions to verified answers, for an operator who had
never seen the task before.

The second lesson comes from the same case study's baseline, and it
reframes what an observability \emph{interface} is for. The raw
captured log proved \emph{sufficient} for complete diagnosis (an
LLM agent given only the data directories answered all three
questions, orphaned delegation included, in under a minute), while
the human operator, given the same files, could not begin. What the
dashboard contributes is therefore the expertise rather than the access:
each view is a frozen form of the queries an
expert would write. Systems built on a sufficient substrate can choose
where that expertise lives (in views for humans, or in agents that
query the log directly), and the shared-log architecture serves both
consumers from the same store.

\section{Future Work}
\label{sec:future}

\paragraph{A second implementation of the contract.} The work that would
most directly test this thesis is a port to another store. The
specification side is done: \S\ref{sec:problem-substrate} states which
features a store must provide and which it may omit at a priced cost,
\S\ref{sec:design-adapters} turns that into the adapter interface, and
\S\ref{sec:design-chronolog} predicts feature by feature what a port would
add and delete. What the contract has not had is an implementation that
could falsify it. Writing two or three (a partitioned commit log such as
Kafka, an embedded analytical store over the archive tiers, a local JSONL
store for single-node development) and checking the predictions of
Table~\ref{tab:substrate-systems} against what each port actually costs
would settle it. A prediction that survives makes the contract a
contribution; one that does not shows where it is under-specified, which is
the more useful outcome.

\paragraph{Upstream convergence.} Two of the three constraints this design
works around are being removed at the source: in ongoing collaboration with
the ChronoLog developers, the blocking replay of C2 and the missing catalog
operation of C3 have since been addressed upstream. Their removal would
retire real machinery here, making the read-only process split a preference
rather than a necessity and the \texttt{\_\_index\_\_} chronicle redundant
for listing. C1, the acceptance window, is inherent to the ingest-optimized
design point and will remain, so the hot/cold split stays; that is the
correct outcome, since it is the part of the read design that is about the
class of substrate rather than one implementation of it. Re-measuring
\S\ref{sec:eval-freshness} against the fixed client would show how much of
Chapter~\ref{chap:design} was design and how much was workaround.

\paragraph{Smaller extensions.} Four follow-ons sit within reach of the
current implementation: ad-hoc queries evaluated along causal edges, the
Pivot-Tracing move~\cite{mace2015pivot}, which the log's ordering and
correlation schema already make expressible; promoting the hot store and
live bus to a small replicated service, so that more than one dashboard can
watch a fleet; learning the diagnoser's rule set from the incident history
rather than hand-building it; and exposing the log and its analyses as an
MCP tool, so that a supervising agent could ask the questions the
case-study operator asked. Live validation at hundreds of real nodes, the
gap \S\ref{sec:eval-limitations} concedes, is chiefly an allocation problem.

\medskip

The direction of the field is not in doubt: agent systems are growing in
both population and physical extent, and the questions operators ask are
increasingly the questions distributed systems have always asked, now
with a conversation attached. This thesis's contribution is to show that
the sixty-year-old toolbox of that discipline, from happened-before to
the shared log, is the one this workload needs.
